\DocumentMetadata{
  lang = pt-BR,
  pdfstandard = A-2b,
  pdfversion = 1.7
}

\documentclass{abntexto-ufc}

\ufcsetup{
  type = scientific-article,
  cover = false,
  catalog-card = false,
  coat-of-arms = false,
  glossary = none,
  index = none,
  author = {Article Body Test Author},
  title = {Article Body Test},
  submission-date = {1 de setembro de 2026},
  approval-date = {5 de setembro de 2026},
  article-author-note = {Universidade Federal do Ceará. article@example.org}
}

\ufcAddBibliographyResource{tests/fixtures/references.bib}

\begin{document}

\ufcPrintArticleFrontMatter{Resumo controlado para o cenário de estrutura e tipografia do corpo do artigo.}

\section{Introdução}

\noindent ARTICLEMARGINCONTROL

ARTICLEBODYSTART Este parágrafo controlado foi escrito com extensão suficiente para ocupar várias linhas naturais no arquivo PDF e permitir a medição física da tipografia do corpo do artigo científico. O texto verifica que o perfil específico mantém corpo nominal de doze pontos, alinhamento justificado e recuo de primeira linha de dois centímetros, ao mesmo tempo em que substitui o espaçamento de uma linha e meia dos trabalhos acadêmicos pelo espaçamento simples exigido para artigos. A continuidade do parágrafo também fornece linhas intermediárias para observar as margens esquerda e direita sem congelar conteúdo editorial como parte do contrato. O cenário permanece deliberadamente simples, reutiliza a infraestrutura compartilhada da classe e termina com um marcador único para delimitar a região medida ARTICLEBODYEND

\par\medskip
\begingroup
\singlesp
\noindent ARTICLECALONE\\
ARTICLECALTWO\\
ARTICLECALTHREE\par
\endgroup

\section{Desenvolvimento}

O desenvolvimento reutiliza a hierarquia de seções compartilhada e contém conteúdo controlado suficiente para comprovar a presença da parte textual obrigatória sem criar uma implementação paralela para artigos.

\section{Considerações finais}

As considerações finais encerram a estrutura textual obrigatória do cenário controlado e preservam a infraestrutura comum de seções e referências.

\nocite{silva2020}
\ufcPrintReferences

\end{document}
